% ===========================================================================
%  Chapitre 7 — Équations différentielles
%  PE 2026, composante ANALYSE
% ===========================================================================
\bmtchapitre{Équations différentielles}
  {Résoudre les équations différentielles du premier ordre et du second ordre,
   avec ou sans second membre, et traiter les cas particuliers
   $y''=m^{2}y$ et $y''=-m^{2}y$.}
  {Analyse \textbullet\ environ 8 heures}

Toutes les équations que vous avez résolues jusqu'ici avaient un nombre pour
inconnue. Celles de ce chapitre ont une \emph{fonction} pour inconnue, et
c'est la seule nouveauté — mais elle change tout. Résoudre, ce n'est plus
trouver un nombre, c'est trouver toutes les fonctions dont la dérivée
entretient un certain rapport avec elles-mêmes.

Ce type d'équation n'est pas une curiosité de fin de programme. C'est le
langage dans lequel s'écrivent la désintégration d'un noyau, le
refroidissement d'un corps, la charge d'un condensateur, l'oscillation d'un
pendule. Chaque fois qu'une grandeur varie proportionnellement à elle-même,
une équation différentielle apparaît — et sa solution contient une
exponentielle. C'est pour cela que ce chapitre vient juste après le
précédent.

% ---------------------------------------------------------------------------
\begin{revision}
Les deux chapitres précédents suffisent.

\begin{enumerate}[label=\textbf{\arabic*.}, leftmargin=*, itemsep=0.35em]
  \item Dériver $x \mapsto e^{3x}$, $x \mapsto e^{-2x}$, $x \mapsto Ke^{ax}$
        où $K$ et $a$ sont des constantes.
  \item Dériver $x \mapsto \cos(2x)$ et $x \mapsto \sin(2x)$, puis calculer
        leurs dérivées secondes.
  \item Résoudre dans $\R$ l'équation $2r^{2}-3r+1 = 0$.
  \item Quel est le discriminant de $r^{2}+4 = 0$ ? Cette équation a-t-elle des
        racines réelles ?
  \item Rappeler la définition d'une primitive.
  \item Si $g'(x) = 0$ sur un intervalle, que peut-on dire de $g$ ?
\end{enumerate}
\end{revision}

\begin{automatismes}
Sans calculatrice, sans brouillon.

\begin{enumerate}[label=\textbf{\alph*)}, leftmargin=*, itemsep=0.25em]
  \item Dérivée de $x \mapsto 5e^{2x}$
  \item Dérivée seconde de $x \mapsto e^{-x}$
  \item Dérivée seconde de $x \mapsto \sin(3x)$
  \item Racines de $r^{2}-5r+6 = 0$
  \item Racines de $r^{2}-4 = 0$
  \item Racines de $r^{2}+9 = 0$ dans $\R$
  \item Dérivée de $x \mapsto Ae^{x}+Be^{-x}$
  \item Une fonction égale à sa propre dérivée
\end{enumerate}
\end{automatismes}

\begin{corrige}[automatismes]
a) $10e^{2x}$ \quad b) $e^{-x}$ \quad c) $-9\sin(3x)$ \quad d) $2$ et $3$ \quad
e) $-2$ et $2$ \quad f) aucune racine réelle \quad
g) $Ae^{x}-Be^{-x}$ \quad h) $x \mapsto e^{x}$, et plus généralement
$x \mapsto Ke^{x}$.
\end{corrige}

% ===========================================================================
\section{Ce qu'on appelle résoudre}

\begin{definition}[équation différentielle, solution]
Une \textbf{équation différentielle} est une égalité reliant une fonction
inconnue $y$, la variable $x$ et une ou plusieurs dérivées de $y$. Une
\textbf{solution} sur un intervalle $I$ est une fonction $\varphi$, dérivable
autant de fois que nécessaire sur $I$, qui vérifie l'égalité pour tout
$x \in I$.

\textbf{Résoudre} l'équation, c'est déterminer \emph{toutes} ses solutions.
\end{definition}

\begin{exemple}[vérifier qu'une fonction est solution]
Montrons que $\varphi(x) = 3e^{-2x}$ est solution de $y' + 2y = 0$ sur $\R$.

On calcule $\varphi'(x) = -6e^{-2x}$, puis
\[
  \varphi'(x) + 2\varphi(x) = -6e^{-2x} + 6e^{-2x} = 0 .
\]
L'égalité est vraie pour tout réel $x$ : $\varphi$ est bien solution.

\smallskip
Notez qu'on n'a rien résolu — on a seulement vérifié. C'est une question de
début d'exercice, et elle se traite toujours ainsi : dériver, remplacer,
constater.
\end{exemple}

\begin{avertissement}
Une équation différentielle n'a jamais une solution unique, mais une famille
entière de solutions, indexée par une ou deux constantes. Écrire « la solution
est $y = 3e^{-2x}$ » est faux tant qu'on n'a pas imposé de condition
initiale : la réponse attendue est « les solutions sont les
$y = Ke^{-2x}$, $K \in \R$ ».
\end{avertissement}

% ===========================================================================
\section{Le premier ordre}

\begin{theoreme}[équation $y' + ay = 0$]
Soit $a$ un réel. Les solutions sur $\R$ de l'équation différentielle
\[
  y' + ay = 0
\]
sont exactement les fonctions
\[
  x \longmapsto K e^{-ax}, \qquad K \in \R .
\]
\end{theoreme}

\begin{demonstration}
\emph{Ces fonctions sont solutions.} Si $\varphi(x) = Ke^{-ax}$, alors
$\varphi'(x) = -aKe^{-ax} = -a\varphi(x)$, donc $\varphi'+a\varphi = 0$.

\smallskip
\emph{Il n'y en a pas d'autres.} Soit $\varphi$ une solution quelconque et
posons $g(x) = \varphi(x)e^{ax}$. Cette fonction est dérivable sur $\R$ et
\[
  g'(x) = \varphi'(x)e^{ax} + a\varphi(x)e^{ax}
        = e^{ax}\underbrace{\bigl[\varphi'(x)+a\varphi(x)\bigr]}_{=\,0} = 0 .
\]
Sur l'intervalle $\R$, une fonction de dérivée nulle est constante : il existe
$K$ tel que $\varphi(x)e^{ax} = K$, c'est-à-dire $\varphi(x) = Ke^{-ax}$.
\end{demonstration}

\bmtmarge{Le procédé de cette preuve — multiplier par $e^{ax}$ pour faire
apparaître une dérivée nulle — s'appelle la méthode du facteur intégrant.
Retenez-en au moins l'idée : on fabrique une fonction dont la dérivée est
nulle.}

\begin{propriete}[condition initiale]
Pour tout couple de réels $(x_{0}, y_{0})$, l'équation $y'+ay = 0$ admet une
\emph{unique} solution vérifiant $y(x_{0}) = y_{0}$ : celle obtenue pour
$K = y_{0}e^{ax_{0}}$.
\end{propriete}

\begin{exemple}[une décroissance radioactive]
La masse $m(t)$ d'un échantillon radioactif vérifie $m' = -\lambda m$, où
$\lambda>0$ est la constante de désintégration. C'est l'équation
$m' + \lambda m = 0$, dont les solutions sont $m(t) = Ke^{-\lambda t}$.

La condition initiale $m(0) = m_{0}$ donne $K = m_{0}$, d'où
\[
  m(t) = m_{0}e^{-\lambda t}.
\]
La demi-vie $T$ est le temps au bout duquel il reste la moitié de la masse :
$e^{-\lambda T} = \tfrac12$, soit $T = \dfrac{\ln 2}{\lambda}$. Elle ne dépend
pas de $m_{0}$ — un gramme et une tonne du même isotope décroissent au même
rythme relatif.
\end{exemple}

\subsection{Avec un second membre}

\begin{theoreme}[structure des solutions]
Considérons l'équation $(E) : y' + ay = h(x)$, où $h$ est une fonction
continue, et l'équation \textbf{homogène associée} $(E_{0}) : y' + ay = 0$.

Si $\varphi_{p}$ est \emph{une} solution particulière de $(E)$, alors les
solutions de $(E)$ sont exactement les fonctions
\[
  \varphi = \varphi_{p} + \varphi_{0},
\]
où $\varphi_{0}$ décrit l'ensemble des solutions de $(E_{0})$.
\end{theoreme}

\begin{demonstration}
Soit $\varphi$ une fonction dérivable. Posons $g = \varphi - \varphi_{p}$.
Alors
\[
  g' + ag = (\varphi'+a\varphi) - (\varphi_{p}'+a\varphi_{p})
          = (\varphi'+a\varphi) - h(x).
\]
Donc $\varphi$ est solution de $(E)$ si et seulement si $g'+ag = 0$,
c'est-à-dire si et seulement si $g$ est solution de $(E_{0})$. Autrement dit
$\varphi = \varphi_{p} + g$ avec $g$ solution de l'homogène.
\end{demonstration}

\begin{methode}{résoudre une équation avec second membre}
La démarche est toujours en trois temps, et les sujets la découpent
explicitement en trois questions.

\begin{enumerate}[label=\textbf{\arabic*.}, leftmargin=*, itemsep=0.28em]
  \item \textbf{Résoudre l'homogène} $(E_{0})$ : on obtient une famille
        dépendant d'une constante.
  \item \textbf{Trouver une solution particulière} de $(E)$. Le sujet la
        fournit presque toujours, en demandant de vérifier qu'une fonction
        donnée convient, ou de déterminer des coefficients pour qu'elle
        convienne.
  \item \textbf{Additionner} : la solution générale est la somme des deux.
\end{enumerate}
\end{methode}

\begin{entrainement}[premier ordre]
\exo[\niveau{1}] Résoudre $y' - 3y = 0$, puis $2y' + y = 0$.

\exo[\niveau{1}] Déterminer la solution de $y' + 5y = 0$ vérifiant $y(0) = 4$.

\exo[\niveau{2}] Vérifier que $\varphi_{p}(x) = 2$ est solution de
$y' + 3y = 6$, puis résoudre cette équation.

\exo[\niveau{3}] Vérifier que $\varphi_{p}(x) = xe^{-x}$ est solution de
$y' + y = e^{-x}$, puis donner la solution générale et celle qui vérifie
$y(0) = 1$.
\end{entrainement}

\begin{corrige}
\textbf{1.} $y = Ke^{3x}$ ; \ pour la seconde, $y' + \tfrac12 y = 0$ donne
$y = Ke^{-x/2}$.

\textbf{2.} $y = Ke^{-5x}$ avec $K = 4$ : $y(x) = 4e^{-5x}$.

\textbf{3.} $\varphi_{p}' + 3\varphi_{p} = 0 + 6 = 6$ : c'est bien une
solution. Les solutions sont donc $y = 2 + Ke^{-3x}$.

\textbf{4.} $\varphi_{p}' = e^{-x} - xe^{-x}$, donc
$\varphi_{p}'+\varphi_{p} = e^{-x}$ : solution particulière. Solution générale
$y = xe^{-x} + Ke^{-x}$. La condition $y(0)=1$ donne $K = 1$, d'où
$y = (x+1)e^{-x}$.
\end{corrige}

% ===========================================================================
\section{Le second ordre}

\begin{definition}[équation caractéristique]
À l'équation différentielle $ay'' + by' + cy = 0$, où $a \neq 0$, on associe
l'équation du second degré
\[
  ar^{2} + br + c = 0,
\]
appelée son \textbf{équation caractéristique}.
\end{definition}

\begin{theoreme}[résolution selon le discriminant]
Soit $\Delta = b^{2}-4ac$ le discriminant de l'équation caractéristique. Les
solutions sur $\R$ de $ay''+by'+cy = 0$ sont :

\medskip
\begin{tabular}{@{}p{2.6cm} p{9.6cm}@{}}
\toprule
\textsf{\footnotesize\bfseries CAS} & \textsf{\footnotesize\bfseries SOLUTIONS} \\
\midrule
$\Delta > 0$ & $y = Ae^{r_{1}x} + Be^{r_{2}x}$, où $r_{1}$ et $r_{2}$ sont les
deux racines réelles distinctes \\
$\Delta = 0$ & $y = (Ax+B)e^{rx}$, où $r$ est la racine double \\
$\Delta < 0$ & $y = e^{\alpha x}\bigl(A\cos\beta x + B\sin\beta x\bigr)$, où
$\alpha \pm i\beta$ sont les racines complexes \\
\bottomrule
\end{tabular}

\medskip
\noindent avec $A$ et $B$ constantes réelles arbitraires.
\end{theoreme}

\begin{remarque}
Deux constantes, et non plus une seule : c'est la marque du second ordre. Il
faut donc \emph{deux} conditions pour les déterminer — typiquement $y(x_{0})$
et $y'(x_{0})$, ou bien deux valeurs de $y$ en deux points.
\end{remarque}

\subsection{Les deux cas particuliers du programme}

\begin{propriete}[$y'' = m^{2}y$ et $y'' = -m^{2}y$]
Soit $m$ un réel strictement positif.
\begin{itemize}[leftmargin=1.4em, itemsep=0.25em]
  \item L'équation $y'' = m^{2}y$, c'est-à-dire $y''-m^{2}y = 0$, a pour
        équation caractéristique $r^{2} = m^{2}$, de racines $m$ et $-m$. Ses
        solutions sont
        \[ y = Ae^{mx} + Be^{-mx}. \]
  \item L'équation $y'' = -m^{2}y$, c'est-à-dire $y''+m^{2}y = 0$, a pour
        équation caractéristique $r^{2} = -m^{2}$, sans racine réelle, de
        racines complexes $\pm im$. Ses solutions sont
        \[ y = A\cos(mx) + B\sin(mx). \]
\end{itemize}
\end{propriete}

\begin{visualisation}[deux signes, deux mondes]
\begin{center}
\begin{tikzpicture}
\begin{axis}[
  width = 11.2cm, height = 5.4cm,
  axis lines = middle, xlabel = {$x$},
  xmin = -0.4, xmax = 6.6, ymin = -2.6, ymax = 3.4,
  xtick = \empty, ytick = \empty,
  axis line style = {bmtGris}, clip = true,
]
  \addplot[bmtLaterite, line width=1.3pt, domain=0:2.4, samples=100] {0.55*exp(0.62*x)};
  \addplot[bmtVetiver, line width=1.3pt, domain=0:6.4, samples=200] {2*cos(deg(1.5*x))};
  \node[bmtLaterite, font=\footnotesize, anchor=west] at (axis cs:2.15,2.9) {$y''=m^{2}y$};
  \node[bmtVetiver, font=\footnotesize, anchor=west] at (axis cs:4.5,1.6) {$y''=-m^{2}y$};
\end{axis}
\end{tikzpicture}
\end{center}

Un seul signe sépare les deux équations, et pourtant leurs solutions n'ont
rien de commun. À gauche, avec le signe $+$, la courbe s'envole : la fonction
et sa dérivée seconde tirent dans le même sens, l'écart s'amplifie. Avec le
signe $-$, la dérivée seconde ramène toujours la fonction vers zéro : elle
oscille indéfiniment sans jamais s'échapper. C'est la différence entre une
réaction en chaîne et un pendule.
\end{visualisation}

\begin{exemple}[les trois cas sur trois équations]
\emph{Cas $\Delta>0$.} $y''-5y'+6y = 0$ : l'équation caractéristique
$r^{2}-5r+6 = 0$ a pour racines $2$ et $3$. Solutions :
$y = Ae^{2x}+Be^{3x}$.

\smallskip
\emph{Cas $\Delta=0$.} $y''-4y'+4y = 0$ : $r^{2}-4r+4 = (r-2)^{2}$, racine
double $2$. Solutions : $y = (Ax+B)e^{2x}$.

\smallskip
\emph{Cas $\Delta<0$.} $y''+2y'+5y = 0$ : $r^{2}+2r+5 = 0$ a pour discriminant
$-16$, donc pour racines $-1 \pm 2i$. Solutions :
$y = e^{-x}\bigl(A\cos 2x + B\sin 2x\bigr)$.
\end{exemple}

\begin{exemple}[avec deux conditions initiales]
Résolvons $y'' + 9y = 0$ avec $y(0) = 2$ et $y'(0) = 3$.

L'équation est du type $y'' = -m^{2}y$ avec $m = 3$ : les solutions sont
$y = A\cos 3x + B\sin 3x$.

\smallskip
La condition $y(0) = 2$ donne $A = 2$. Puis
$y'(x) = -3A\sin 3x + 3B\cos 3x$, donc $y'(0) = 3B = 3$ et $B = 1$.

\smallskip
La solution cherchée est $y(x) = 2\cos 3x + \sin 3x$, et c'est la seule.
\end{exemple}

\begin{entrainement}[second ordre]
\exo[\niveau{1}] Résoudre $y''-y = 0$, puis $y''+4y = 0$.

\exo[\niveau{2}] Résoudre $y''-6y'+9y = 0$, puis $y''+y'-2y = 0$.

\exo[\niveau{2}] Résoudre $y''+2y'+2y = 0$.

\exo[\niveau{3}] Déterminer la solution de $y''-4y = 0$ vérifiant $y(0) = 1$
et $y'(0) = 0$, puis montrer qu'elle est paire.
\end{entrainement}

\begin{corrige}
\textbf{1.} $r^{2}-1 = 0$ : $y = Ae^{x}+Be^{-x}$. \ $r^{2}+4 = 0$ :
$y = A\cos 2x + B\sin 2x$.

\textbf{2.} $(r-3)^{2} = 0$ : $y = (Ax+B)e^{3x}$. \ $r^{2}+r-2 = 0$ de racines
$1$ et $-2$ : $y = Ae^{x}+Be^{-2x}$.

\textbf{3.} $\Delta = -4$, racines $-1\pm i$ :
$y = e^{-x}(A\cos x + B\sin x)$.

\textbf{4.} $y = Ae^{2x}+Be^{-2x}$ ; $y(0) = A+B = 1$ et
$y'(0) = 2A-2B = 0$ donnent $A = B = \tfrac12$, soit
$y = \dfrac{e^{2x}+e^{-2x}}{2}$. On vérifie $y(-x) = y(x)$ : la fonction est
paire.
\end{corrige}

% ===========================================================================
% ===========================================================================
\begin{annales}
Les équations différentielles apparaissent au baccalauréat sous deux visages :
une question courte et technique, ou bien une partie entière du problème
d'analyse, articulée à l'étude d'une fonction. Les énoncés qui suivent sont
repris de sessions réelles ; leur provenance est indiquée à chaque fois.

\exoann{Baccalauréat, Madagascar, série C, session 2012 — problème 2, partie B}
On considère l'équation différentielle $(E_{1}) : y'-y = xe^{x}$.
\emph{Le second membre étant $xe^{x}$, on cherche une solution particulière
sous la forme $u(x) = \left(ax^{2}+bx\right)e^{x}$.}
\begin{enumerate}[label=\textbf{\arabic*)}, leftmargin=*, itemsep=0.15em]
  \item Résoudre l'équation différentielle $(E_{2}) : y'-y = 0$.
  \item Déterminer $a$ et $b$ pour que $u$ soit solution de $(E_{1})$.
  \item $\varphi$ étant dérivable sur $\R$, montrer que $\varphi+u$ est
        solution de $(E_{1})$ si et seulement si $\varphi$ est solution de
        $(E_{2})$. En déduire l'ensemble des solutions de $(E_{1})$.
  \item Déterminer la solution de $(E_{1})$ qui s'annule en $0$.
\end{enumerate}

\exoann{Baccalauréat, Cameroun, séries D et TI, session 2024 — exercice 3}
On considère $(E) : y''+4y'+4y = 0$ et $(E') : y''+4y'+4y = -4$. Résoudre
$(E)$, puis $(E')$.

\exoann{Baccalauréat, Cameroun, séries D et TI, session 2024 — exercice 3, suite}
Soit $g(x) = -1-(x+0{,}5)e^{-2x}$. Démontrer que $g$ est la solution de
$(E')$ vérifiant $g(0) = -1{,}5$ et $g'(0) = 0$.

\exoann{Baccalauréat, Madagascar, série S, session 2026 — problème 2, partie B}
On considère $(E_{1}) : 2y''-3y'+y = (ax+b)e^{-x}$ et
$(E_{2}) : 2y''-3y'+y = 0$, où $a$ et $b$ sont deux réels.
\begin{enumerate}[label=\textbf{\arabic*)}, leftmargin=*, itemsep=0.15em]
  \item Déterminer $a$ et $b$ pour que $f_{1}(x) = (x+1)e^{-x}$ soit solution
        de $(E_{1})$.
  \item Montrer que $f_{1}+g$ est solution de $(E_{1})$ si et seulement si $g$
        est solution de $(E_{2})$.
  \item En déduire l'ensemble des solutions de $(E_{1})$.
\end{enumerate}

\exoann{Baccalauréat, Madagascar, série S, session 2023 — problème 2}
On considère l'équation différentielle
$(E) : y'-2y = \dfrac{-2}{1+e^{-2x}}$.
\begin{enumerate}[label=\textbf{\alph*)}, leftmargin=*, itemsep=0.15em]
  \item Vérifier que $f_{1}(x) = e^{2x}\ln\!\left(1+e^{-2x}\right)$ est
        solution de $(E)$.
  \item Résoudre $(E_{0}) : y'-2y = 0$, puis en déduire l'ensemble des
        solutions de $(E)$.
\end{enumerate}

\exoann{Baccalauréat, Côte d'Ivoire, série C, session 2023 — exercice 6}
La teneur en carbone 14 d'un organisme mort décroît selon
$P'(t) = -\lambda P(t)$, où $t$ est exprimé en années,
$\lambda = 1{,}2444\times10^{-4}$ et $P(0) = P_{0}$.
\begin{enumerate}[label=\textbf{\alph*)}, leftmargin=*, itemsep=0.15em]
  \item Résoudre cette équation différentielle et exprimer $P(t)$ en fonction
        de $P_{0}$, $\lambda$ et $t$.
  \item Un fragment d'os présente $70\,\%$ de la teneur initiale. Déterminer
        son âge, arrondi à l'année.
\end{enumerate}

\exoann{Baccalauréat, Côte d'Ivoire, série C, session 2024 — exercice 2}
Résoudre dans $\R$ l'équation différentielle $y'+2y = 0$, puis déterminer la
solution vérifiant $y(0) = 3$.

\exoann{D'après le baccalauréat, Cameroun, séries C et E, session 2024 — exercice 1}
Une urne contient six boules numérotées de $1$ à $6$. On tire successivement
et avec remise deux boules ; on note $a$ le premier numéro et $b$ le second.
On considère alors $(E) : y''+ay'+by = 0$, d'équation caractéristique
$r^{2}+ar+b = 0$.
\begin{enumerate}[label=\textbf{\alph*)}, leftmargin=*, itemsep=0.15em]
  \item Déterminer le nombre de couples $(a\,;b)$ pour lesquels $(E)$ admet
        deux solutions réelles, distinctes ou confondues.
  \item En déduire le nombre de couples pour lesquels l'équation
        caractéristique n'a pas de racine réelle.
\end{enumerate}
\end{annales}

\begin{corrige}[annales]
\textbf{1.} Les solutions de $(E_{2})$ sont les $x\mapsto Ce^{x}$, $C$ réel.
Avec $u(x) = \left(ax^{2}+bx\right)e^{x}$, on obtient
$u'-u = (2ax+b)e^{x}$ ; l'identification avec $xe^{x}$ donne $a = \frac12$ et
$b = 0$, soit $u(x) = \frac{x^{2}}{2}e^{x}$. Ensuite
$(\varphi+u)'-(\varphi+u) = \left(\varphi'-\varphi\right)+\left(u'-u\right)
= \left(\varphi'-\varphi\right)+xe^{x}$ : cette expression vaut $xe^{x}$ si et
seulement si $\varphi'-\varphi = 0$. Les solutions de $(E_{1})$ sont donc les
$x\mapsto\left(\frac{x^{2}}{2}+C\right)e^{x}$. Celle qui s'annule en $0$
correspond à $C = 0$ : $y(x) = \frac{x^{2}}{2}e^{x}$.

\textbf{2.} L'équation caractéristique $r^{2}+4r+4 = 0$ a la racine double
$r = -2$ : les solutions de $(E)$ sont les $x\mapsto(Ax+B)e^{-2x}$. La
fonction constante $y = -1$ est solution particulière de $(E')$, dont les
solutions sont donc les $x\mapsto(Ax+B)e^{-2x}-1$.

\textbf{3.} $g$ correspond à $A = -1$ et $B = -0{,}5$ : c'est bien une
solution de $(E')$. De plus $g(0) = -1-0{,}5 = -1{,}5$ et
$g'(x) = -e^{-2x}+2(x+0{,}5)e^{-2x} = 2xe^{-2x}$, donc $g'(0) = 0$. L'unicité
de la solution d'un tel problème assure que c'est la seule.

\textbf{4.} Avec $f_{1}(x) = (x+1)e^{-x}$ on trouve $f_{1}'(x) = -xe^{-x}$ et
$f_{1}''(x) = (x-1)e^{-x}$, d'où
$2f_{1}''-3f_{1}'+f_{1} = \left[2(x-1)+3x+(x+1)\right]e^{-x} = (6x-1)e^{-x}$ :
donc $a = 6$ et $b = -1$. La linéarité donne aussitôt l'équivalence
demandée. Enfin $2r^{2}-3r+1 = 0$ a pour racines $1$ et $\frac12$ : les
solutions de $(E_{1})$ sont les
$x\mapsto(x+1)e^{-x}+Ae^{x}+Be^{x/2}$.

\textbf{5. a)} En dérivant,
$f_{1}'(x) = 2e^{2x}\ln\!\left(1+e^{-2x}\right)
-\dfrac{2e^{2x}e^{-2x}}{1+e^{-2x}} = 2f_{1}(x)-\dfrac{2}{1+e^{-2x}}$, ce qui
est exactement $(E)$.
\textbf{b)} Les solutions de $(E_{0})$ sont les $x\mapsto Ce^{2x}$ ; celles de
$(E)$ sont donc les $x\mapsto e^{2x}\ln\!\left(1+e^{-2x}\right)+Ce^{2x}$.

\textbf{6. a)} L'équation $P' = -\lambda P$ donne
$P(t) = P_{0}e^{-\lambda t}$.
\textbf{b)} $e^{-\lambda t} = 0{,}7$ donne
$t = \dfrac{-\ln 0{,}7}{\lambda} = \dfrac{0{,}356675}{1{,}2444\times10^{-4}}
\approx 2866$ ans.

\textbf{7.} Les solutions sont les $x\mapsto Ce^{-2x}$ ; la condition
$y(0) = 3$ impose $C = 3$, soit $y(x) = 3e^{-2x}$.

\textbf{8. a)} Il faut $\Delta = a^{2}-4b \geq 0$, soit $b \leq
\frac{a^{2}}{4}$. Pour $a = 1$ aucun $b$ ne convient ; pour $a = 2$, un seul ;
pour $a = 3$, deux ; pour $a = 4$, quatre ; pour $a = 5$ et pour $a = 6$, les
six valeurs conviennent. Au total $0+1+2+4+6+6 = 19$ couples.
\textbf{b)} Il y a $36$ couples possibles, donc $36-19 = 17$ couples pour
lesquels l'équation caractéristique n'a pas de racine réelle.
\end{corrige}

\begin{jecherche}[le refroidissement d'un four à briques]
Un four artisanal est éteint à l'instant $t = 0$. Sa température $\theta(t)$,
en degrés Celsius, $t$ heures plus tard, suit la loi de Newton : la vitesse de
refroidissement est proportionnelle à l'écart entre la température du four et
celle de l'air ambiant, constante et égale à $25\,^{\circ}$C. On écrit donc
\[
  \theta'(t) = -k\bigl(\theta(t) - 25\bigr), \qquad k>0 .
\]
À l'extinction, le four est à $625\,^{\circ}$C ; deux heures plus tard, il est
à $325\,^{\circ}$C.

\begin{enumerate}[label=\textbf{\arabic*.}, leftmargin=*, itemsep=0.3em]
  \item On pose $u(t) = \theta(t)-25$. Montrer que $u$ vérifie une équation
        différentielle du premier ordre sans second membre, et la résoudre.
  \item En déduire l'expression de $\theta(t)$ en fonction de $k$.
  \item Utiliser les deux mesures pour déterminer $k$, puis donner
        l'expression complète de $\theta$.
  \item Au bout de combien de temps le four descend-il sous $100\,^{\circ}$C ?
  \item Le four peut-il atteindre la température ambiante ? Que dit le modèle,
        et qu'en pensez-vous ?
  \item L'atelier voisin, moins isolé, a une constante $k$ deux fois plus
        grande. Sans refaire tout le calcul, comparer les deux durées de
        refroidissement jusqu'à $100\,^{\circ}$C.
\end{enumerate}
\end{jecherche}

\begin{corrige}[le refroidissement d'un four à briques]
\textbf{1.} Avec $u = \theta-25$ on a $u' = \theta'$, donc
$u' = -k\,u$, c'est-à-dire $u'+ku = 0$. Les solutions sont
$u(t) = Ke^{-kt}$.

\textbf{2.} Donc $\theta(t) = 25 + Ke^{-kt}$. La condition $\theta(0) = 625$
donne $K = 600$ :
$\theta(t) = 25 + 600e^{-kt}$.

\textbf{3.} $\theta(2) = 325$ donne $600e^{-2k} = 300$, soit
$e^{-2k} = \tfrac12$ et $k = \dfrac{\ln 2}{2} \approx 0{,}347$. Ainsi
\[
  \theta(t) = 25 + 600 \times 2^{-t/2}.
\]

\textbf{4.} $\theta(t) < 100$ équivaut à $600 \times 2^{-t/2} < 75$, soit
$2^{-t/2} < 2^{-3}$ et donc $t > 6$. Le four passe sous cent degrés après six
heures.

\textbf{5.} $\lim_{t\to+\infty}\theta(t) = 25$ : la température ambiante n'est
atteinte qu'à la limite, jamais en temps fini. C'est une propriété du modèle,
non du four — en pratique, l'écart devient indétectable au bout de quelques
heures.

\textbf{6.} Le temps cherché vérifie $600e^{-kt} = 75$, soit
$t = \dfrac{3\ln 2}{k}$ : il est inversement proportionnel à $k$. Un atelier
dont la constante est deux fois plus grande refroidit donc deux fois plus
vite, et atteint cent degrés en trois heures.
\end{corrige}


% ===========================================================================
\begin{jemeteste}
Répondre par vrai ou faux, en justifiant en une ligne.

\begin{enumerate}[label=\textbf{\arabic*.}, leftmargin=*, itemsep=0.28em]
  \item L'équation $y'+2y = 0$ admet une unique solution.
  \item Les solutions de $y'-3y=0$ sont les $x \mapsto Ke^{-3x}$.
  \item Une équation du second ordre a des solutions dépendant de deux
        constantes.
  \item Si l'équation caractéristique n'a pas de racine réelle, l'équation
        différentielle n'a pas de solution.
  \item Les solutions de $y''+16y = 0$ sont périodiques.
  \item Pour résoudre une équation avec second membre, il suffit de connaître
        une solution particulière et l'équation homogène associée.
\end{enumerate}
\end{jemeteste}

\begin{corrige}[je me teste]
\textbf{1.} Faux — elle en admet une infinité, une par valeur de $K$ ; l'unicité
n'apparaît qu'avec une condition initiale.
\textbf{2.} Faux — ce sont les $Ke^{3x}$ : attention au signe.
\textbf{3.} Vrai.
\textbf{4.} Faux — les solutions font alors intervenir cosinus et sinus.
\textbf{5.} Vrai — ce sont les $A\cos 4x + B\sin 4x$, de période
$\dfrac{\pi}{2}$.
\textbf{6.} Vrai — c'est exactement le théorème de structure.
\end{corrige}

% ===========================================================================
\begin{commeaubac}
\bmtsource{Baccalauréat, session 2026 — série S, problème 2, partie B}
On considère les deux équations différentielles
\[
  (E_{1}) : 2y''-3y'+y = (ax+b)e^{-x},
  \qquad
  (E_{2}) : 2y''-3y'+y = 0,
\]
où $a$ et $b$ sont deux réels.

\begin{enumerate}[label=\textbf{\arabic*.}, leftmargin=*, itemsep=0.25em]
  \item Déterminer les réels $a$ et $b$ pour que $f_{1}(x) = (x+1)e^{-x}$ soit
        une solution de $(E_{1})$ dans $\R$. \hfill\textup{(0,5 pt)}
  \item Soit $g$ une fonction deux fois dérivable sur $\R$.
        \begin{enumerate}[label=\textbf{\alph*)}, leftmargin=1.4em]
          \item Montrer que $f_{1}+g$ est solution de $(E_{1})$ si et seulement
                si $g$ est solution de $(E_{2})$. \hfill\textup{(0,5 pt)}
          \item En déduire la solution générale de $(E_{1})$ dans $\R$.
                \hfill\textup{(0,75 pt)}
        \end{enumerate}
\end{enumerate}

\medskip
\textbf{Partie C.} Pour tout $n \in \N^{*}$, on pose
$\displaystyle I_{n} = \int_{-1}^{0}(x+1)^{n}e^{-x}\dx$.

\begin{enumerate}[label=\textbf{\arabic*.}, leftmargin=*, itemsep=0.25em, start=3]
  \item Calculer $I_{1}$ et interpréter graphiquement ce résultat.
        \hfill\textup{(0,75 pt)}
  \item Étudier la variation de la suite $(I_{n})$, puis en déduire qu'elle est
        convergente. \hfill\textup{(1 pt)}
  \item \begin{enumerate}[label=\textbf{\alph*)}, leftmargin=1.4em]
          \item Montrer que pour tout $n \in \N^{*}$ et pour $-1<x<0$,
                \[ \frac{1}{n+1} \leq I_{n} \leq \frac{e}{n+1}. \]
                \hfill\textup{(0,75 pt)}
          \item En déduire $\displaystyle\lim_{n \to +\infty} I_{n}$.
                \hfill\textup{(0,25 pt)}
        \end{enumerate}
\end{enumerate}
\end{commeaubac}

\begin{corrige}[baccalauréat 2026, série S]
\textbf{1.} $f_{1}(x) = (x+1)e^{-x}$ donne $f_{1}'(x) = -xe^{-x}$ puis
$f_{1}''(x) = (x-1)e^{-x}$. Alors
\[
  2f_{1}''-3f_{1}'+f_{1}
  = \bigl[2(x-1)+3x+(x+1)\bigr]e^{-x} = (6x-1)e^{-x}.
\]
Donc $a = 6$ et $b = -1$.

\textbf{2. a)} Par linéarité, $2(f_{1}+g)''-3(f_{1}+g)'+(f_{1}+g)
= \underbrace{\left[2f_{1}''-3f_{1}'+f_{1}\right]}_{=(ax+b)e^{-x}}
+ \left[2g''-3g'+g\right]$. Cette somme vaut $(ax+b)e^{-x}$ si et seulement si
le second crochet est nul, c'est-à-dire si $g$ est solution de $(E_{2})$.

\textbf{b)} L'équation caractéristique de $(E_{2})$ est $2r^{2}-3r+1 = 0$, de
discriminant $1$ et de racines $1$ et $\tfrac12$. Donc
$g(x) = Ae^{x}+Be^{x/2}$, et la solution générale de $(E_{1})$ est
\[
  y(x) = (x+1)e^{-x} + Ae^{x} + Be^{\frac{x}{2}}, \qquad (A,B) \in \R^{2}.
\]

\textbf{3.} Intégration par parties avec $u = x+1$ et $v' = e^{-x}$ :
\[
  I_{1} = \Bigl[-(x+1)e^{-x}\Bigr]_{-1}^{0} + \int_{-1}^{0}e^{-x}\dx
        = -1 + \Bigl[-e^{-x}\Bigr]_{-1}^{0} = -1 + (e-1) = e-2 .
\]
Comme $f_{1} \geq 0$ sur $\intervalle{-1}{0}$, ce nombre est l'aire, en unités
d'aire, du domaine compris entre la courbe de $f_{1}$, l'axe des abscisses et
les droites $x=-1$ et $x=0$.

\textbf{4.} $I_{n+1}-I_{n} = \displaystyle\int_{-1}^{0}(x+1)^{n}\bigl[(x+1)-1\bigr]e^{-x}\dx
= \int_{-1}^{0}x(x+1)^{n}e^{-x}\dx \leq 0$, car sur $\intervalle{-1}{0}$ on a
$x \leq 0$ et $(x+1)^{n}e^{-x} \geq 0$. La suite est donc décroissante ; elle
est minorée par $0$ puisque l'intégrande est positif. Décroissante et minorée,
elle converge.

\textbf{5. a)} Sur $\intervalle{-1}{0}$, $0 \leq -x \leq 1$ donc
$1 \leq e^{-x} \leq e$. En multipliant par $(x+1)^{n} \geq 0$ puis en
intégrant :
\[
  \int_{-1}^{0}(x+1)^{n}\dx \ \leq\ I_{n} \ \leq\ e\int_{-1}^{0}(x+1)^{n}\dx .
\]
Or $\displaystyle\int_{-1}^{0}(x+1)^{n}\dx = \left[\frac{(x+1)^{n+1}}{n+1}\right]_{-1}^{0}
= \frac{1}{n+1}$, d'où l'encadrement annoncé.

\textbf{b)} Les deux bornes tendent vers $0$ : par le théorème des gendarmes,
$\lim I_{n} = 0$.
\end{corrige}

\begin{commeaubac}
\bmtsource{Exercice au format de l'épreuve}
On considère l'équation différentielle
$(E) : y'' - 4y' + 4y = 2e^{2x}$.

\begin{enumerate}[label=\textbf{\arabic*.}, leftmargin=*, itemsep=0.25em]
  \item Résoudre l'équation homogène associée $(E_{0}) : y''-4y'+4y = 0$.
        \hfill\textup{(1 pt)}
  \item Déterminer le réel $\lambda$ pour que $\varphi(x) = \lambda x^{2}e^{2x}$
        soit solution de $(E)$. \hfill\textup{(1,5 pt)}
  \item En déduire la solution générale de $(E)$. \hfill\textup{(0,75 pt)}
  \item Déterminer la solution de $(E)$ vérifiant $y(0) = 1$ et $y'(0) = 0$.
        \hfill\textup{(1,25 pt)}
  \item Soit $h$ la solution obtenue à la question précédente. Étudier le signe
        de $h$ sur $\R$ et déterminer $\limminf h(x)$.
        \hfill\textup{(1,5 pt)}
\end{enumerate}
\end{commeaubac}

\begin{corrige}[exercice au format de l'épreuve]
\textbf{1.} $r^{2}-4r+4 = (r-2)^{2}$ : racine double $2$, donc
$y = (Ax+B)e^{2x}$.

\textbf{2.} $\varphi(x) = \lambda x^{2}e^{2x}$ donne
$\varphi'(x) = \lambda(2x+2x^{2})e^{2x}$ et
$\varphi''(x) = \lambda(2+8x+4x^{2})e^{2x}$. Alors
\[
  \varphi''-4\varphi'+4\varphi
  = \lambda\bigl[(2+8x+4x^{2}) - 4(2x+2x^{2}) + 4x^{2}\bigr]e^{2x}
  = 2\lambda e^{2x}.
\]
L'égalité avec $2e^{2x}$ impose $\lambda = 1$.

\textbf{3.} $y(x) = x^{2}e^{2x} + (Ax+B)e^{2x} = \left(x^{2}+Ax+B\right)e^{2x}$.

\textbf{4.} $y(0) = B = 1$. Puis
$y'(x) = \left(2x+A\right)e^{2x} + 2\left(x^{2}+Ax+B\right)e^{2x}$, donc
$y'(0) = A + 2B = 0$, soit $A = -2$. La solution est
$h(x) = \left(x^{2}-2x+1\right)e^{2x} = (x-1)^{2}e^{2x}$.

\textbf{5.} $(x-1)^{2} \geq 0$ et $e^{2x}>0$ : $h \geq 0$ sur $\R$, avec
$h(x) = 0$ uniquement en $x=1$. En $-\infty$, $e^{2x} \to 0$ l'emporte sur la
croissance polynomiale : $\limminf h(x) = 0$.
\end{corrige}

% ===========================================================================
\begin{concours}
\bmtsource{Vers les classes préparatoires}
On considère l'équation différentielle $(E) : y'' + y = \cos x$.

\begin{enumerate}[label=\textbf{\arabic*.}, leftmargin=*, itemsep=0.25em]
  \item Résoudre l'équation homogène associée.
  \item Montrer qu'aucune fonction de la forme $x \mapsto \lambda\cos x$ ne
        peut être solution de $(E)$.
  \item Déterminer un réel $\mu$ tel que
        $\varphi(x) = \mu x\sin x$ soit solution de $(E)$.
  \item En déduire la solution générale de $(E)$, puis celle qui vérifie
        $y(0) = 0$ et $y'(0) = 0$.
  \item La solution obtenue est-elle bornée sur $\R$ ? Commenter le phénomène.
\end{enumerate}
\end{concours}

\begin{corrige}[vers le concours]
\textbf{1.} $r^{2}+1 = 0$ n'a pas de racine réelle, ses racines complexes sont
$\pm i$ : les solutions de l'homogène sont
$y = A\cos x + B\sin x$.

\textbf{2.} Si $\varphi(x) = \lambda\cos x$, alors
$\varphi''(x) = -\lambda\cos x$ et $\varphi''+\varphi = 0$, quel que soit
$\lambda$. Une telle fonction ne peut donc jamais produire le second membre
$\cos x$ : elle est déjà solution de l'homogène.

\textbf{3.} Pour $\varphi(x) = \mu x\sin x$ :
$\varphi'(x) = \mu(\sin x + x\cos x)$ puis
$\varphi''(x) = \mu(2\cos x - x\sin x)$. Donc
\[
  \varphi''+\varphi = \mu\bigl(2\cos x - x\sin x\bigr) + \mu x\sin x
                    = 2\mu\cos x .
\]
L'égalité avec $\cos x$ impose $\mu = \dfrac12$.

\textbf{4.} La solution générale est
$y = \dfrac{x\sin x}{2} + A\cos x + B\sin x$. La condition $y(0) = 0$ donne
$A = 0$ ; en dérivant, $y'(0) = B = 0$. La solution cherchée est donc
$y(x) = \dfrac{x\sin x}{2}$.

\textbf{5.} Elle n'est pas bornée : pour $x = \dfrac{\pi}{2}+2k\pi$, on a
$\sin x = 1$ et $y(x) = \dfrac{x}{2} \to +\infty$. Le second membre oscille
exactement à la fréquence propre du système, et l'amplitude croît sans limite :
c'est la résonance. Un pont excité à sa fréquence propre subit le même
phénomène.
\end{corrige}


% =========================================================================
\begin{devoir}[2 heures]
Trois exercices indépendants, notés sur $20$. Le devoir est cumulatif : il porte sur ce chapitre et sur ceux qui le précèdent. Les énoncés sont repris de sessions réelles et assemblés ici en épreuve ; leur provenance est indiquée à chaque fois.

\exods{1}{Baccalauréat, Congo-Brazzaville, série C, session 2023 — exercice 3}{5 points}
Soit $f(x) = \dfrac{6\ln(x+3)}{x+3}$ sur $\intervallefo{0}{+\infty}$.
Démontrer que $f'(x) = \dfrac{6\left[1-\ln(x+3)\right]}{(x+3)^{2}}$, en
déduire le sens de variation de $f$, puis calculer
$\lim_{x\to+\infty}f(x)$ et dresser le tableau de variation.

\exods{2}{Baccalauréat, Côte d'Ivoire, série C, session 2024 — exercice 5}{7 points}
Soit $f(x) = e^{\sqrt x}$ sur $\intervallefo{0}{+\infty}$.
\begin{enumerate}[label=\textbf{\alph*)}, leftmargin=*, itemsep=0.15em]
  \item Calculer $\lim_{x\to+\infty}f(x)$.
  \item Étudier la dérivabilité de $f$ en $0$ et interpréter graphiquement.
  \item Calculer $f'(x)$ pour $x>0$, dresser le tableau de variation et
        construire la courbe.
\end{enumerate}

\exods{3}{Baccalauréat, Côte d'Ivoire, série C, session 2023 — exercice 6}{8 points}
La teneur en carbone 14 d'un organisme mort décroît selon
$P'(t) = -\lambda P(t)$, où $t$ est exprimé en années,
$\lambda = 1{,}2444\times10^{-4}$ et $P(0) = P_{0}$.
\begin{enumerate}[label=\textbf{\alph*)}, leftmargin=*, itemsep=0.15em]
  \item Résoudre cette équation différentielle et exprimer $P(t)$ en fonction
        de $P_{0}$, $\lambda$ et $t$.
  \item Un fragment d'os présente $70\,\%$ de la teneur initiale. Déterminer
        son âge, arrondi à l'année.
\end{enumerate}

\end{devoir}

\begin{corrige}[devoir surveillé]
\textbf{Exercice 1.} $f'(x) = 6\,\dfrac{\frac{1}{x+3}(x+3)-\ln(x+3)}{(x+3)^{2}}
= \dfrac{6\left[1-\ln(x+3)\right]}{(x+3)^{2}}$. Pour $x \geq 0$ on a
$x+3 \geq 3 > e$, donc $\ln(x+3)>1$ et $f'(x)<0$ : $f$ est strictement
décroissante. Comme $\frac{\ln X}{X}\to0$, $f(x)\to0$ en $+\infty$ ; $f$
décroît donc de $f(0) = 2\ln3$ à $0$.

\textbf{Exercice 2. a)} $\sqrt x\to+\infty$ donc $f(x)\to+\infty$.
\textbf{b)} Le taux vaut $\dfrac{e^{\sqrt x}-1}{x}
= \dfrac{e^{\sqrt x}-1}{\sqrt x}\cdot\dfrac{1}{\sqrt x}$ ; le premier facteur
tend vers $1$ et le second vers $+\infty$ : la limite est $+\infty$, $f$ n'est
pas dérivable en $0$ et la courbe y admet une demi-tangente verticale.
\textbf{c)} Pour $x>0$, $f'(x) = \dfrac{e^{\sqrt x}}{2\sqrt x}>0$ : $f$ croît
de $1$ à $+\infty$.

\textbf{Exercice 3. a)} L'équation $P' = -\lambda P$ donne
$P(t) = P_{0}e^{-\lambda t}$.
\textbf{b)} $e^{-\lambda t} = 0{,}7$ donne
$t = \dfrac{-\ln 0{,}7}{\lambda} = \dfrac{0{,}356675}{1{,}2444\times10^{-4}}
\approx 2866$ ans.

\end{corrige}
